%
% Partialbruchzeregung des gebrochenen Teils
%
\subsection{Partialbruchzerlegung des gebrochenen Teils}
Es bleibt jetzt also noch, das Integral der Form
\[
\int \frac{w(x)}{\vartheta}\,dx
\]
für eine rationale Funktion $w(x)\in\mathbb{Q}(x)$ zu berechnen.
Die rationale Funktion $w(x)$ hat eine Partialbruchzerlegung
\begin{align*}
w(z)
&=
\sum_{i=1}^n \sum_{k=1}^{n_i}\frac{A_{ki}}{(x-\xi_i)^k},
\end{align*}
wobei $\xi_i$ die Nullstellen des Nenners mit Vielfachheit $n_k$ sind.
Das Integral wird dann zu
\begin{align*}
\int \frac{w(x)}{\vartheta}\,dx
&=
\sum_{i=1}^n \sum_{k=1}^{n_i}
A_{ki}
\int
\frac{1}{(x-\xi_i)^k\vartheta}\,dx.
\end{align*}
Der gebrochene Teil kann also mit der Partialbruchzerlegung auf die
Integrale 
\begin{equation}
\int\frac{1}{\vartheta}\,dx
\qquad\text{und}\qquad
\int\frac{1}{(x-\xi)^k\vartheta}\,dx
\label{buch:intalg:wurzel:eqn:grundintegrale}
\end{equation}
zurückgeführt werden.
Zusätzlich zum Integral mit Integrand $1/\vartheta$ treten jetzt noch
weitere Integrale auf.
